\section{Deninger's determinant package}

The oscillator and prime diagonal give the scalar identity of
Corollary~\ref{cor:completed-euler}. An arithmetic cohomological
realization must produce its operators independently, with the same
normalization. The determinant prescription extends to the following
algebraic spectral setting; the required continuation is an explicit
hypothesis.

\begin{definition}[Zeta-regularized determinant]
\label{def:det}
Let \(V=\bigoplus_\alpha V_\alpha\) be a countable algebraic direct sum of
finite-dimensional \(T\)-invariant complex vector spaces.  Eigenvalues are
counted with algebraic multiplicity.  Use
\[
 \lambda^{-z}=\exp\!\bigl(-z(\log|\lambda|+i\arg\lambda)\bigr),
 \qquad -\pi<\arg\lambda\leq\pi.
\]
If \(0\) is an eigenvalue, set \(\det_\infty(T\mid V)=0\).  Otherwise
suppose
\[
 \zeta_T(z)=\sum_{\lambda\in\operatorname{Spec}_{\mathrm{alg}}T}
 m_T(\lambda)\lambda^{-z}
\]
converges absolutely in a right half-plane, extends meromorphically to
\(\operatorname{Re}z>-\varepsilon\) for some \(\varepsilon>0\), and is
holomorphic at zero.  Then
\[
 \det_\infty(T\mid V)=\exp(-\zeta_T'(0)).
\]
For a real space, this definition is applied after complexification.
\end{definition}

When zero is not an eigenvalue and the spectral zeta hypotheses hold,
for \(c>0\) compatible branches give
\[
 \det_\infty(cT)=c^{\zeta_T(0)}\det_\infty(T).
\]
Thus a factor such as \(2\pi\) has a scale anomaly and cannot be suppressed
without computing \(\zeta_T(0)\).

\begin{theorem}[Conditional Deninger determinant formula]
\label{thm:deninger}
Assume that real vector spaces \(H^i\), concentrated in degrees
\(i=0,1,2\), and endomorphisms \(\Theta_i\) are given, and that each
complexification is a countable algebraic direct sum of finite-dimensional
\(\Theta_i\)-invariant subspaces as in Definition~\ref{def:det}.  For
\(s\notin\operatorname{Spec}_{\mathrm{alg}}(\Theta_i)\), apply
Definition~\ref{def:det} to the pencil
\[
 D_i(s)=\det_\infty\left(
 \frac{s-\Theta_i}{2\pi}\,\middle|\,H^i\right).
\]
Assume that each locally defined determinant is the restriction of a
single-valued meromorphic function \(D_i\) on \(\mathbb C\), with the zero
value of Definition~\ref{def:det} at spectral parameters, and that
\[
 (H^0,\Theta_0)=(\mathbb R,0),
 \qquad (H^2,\Theta_2)=(\mathbb R,1),
\]
and that Deninger's cohomological identity holds in his normalization:
\[
 \widehat\zeta_D(s)=\frac{D_1(s)}{D_0(s)D_2(s)},
 \qquad \widehat\zeta_D(s)=2^{-1/2}\Lambda(s).
\]
Then
\[
 D_0(s)=\frac{s}{2\pi},\qquad
 D_2(s)=\frac{s-1}{2\pi},
\]
and
\[
 \Lambda(s)=\sqrt2\,\frac{D_1(s)}{D_0(s)D_2(s)},
 \qquad
 \xi_{\mathbb R}(s)=2\sqrt2\,\pi^2D_1(s).
\]
\end{theorem}

\begin{proof}
The determinants in degrees zero and two are ordinary one-dimensional
determinants.  Substitution into the assumed alternating determinant formula
gives
\[
 D_1(s)=2^{-1/2}\Lambda(s)\frac{s(s-1)}{4\pi^2}
 =\frac{\xi_{\mathbb R}(s)}{2\sqrt2\,\pi^2}.
\]
The assumed factor $2^{-1/2}$ is the normalization in
\cite[\S3, equation~(3)]{Deninger1998}, and it agrees with
Corollary~\ref{cor:completed-euler}.
\end{proof}

The theorem is an implication from the displayed package.  It does not
construct \(H^1\), \(\Theta_1\), or their determinant.  In particular,
\(\det_\infty(\Theta_1)\) is a scalar, whereas
\(\xi_{\mathbb R}(s)\) is a function; the parameter belongs in the pencil
\((s-\Theta_1)/(2\pi)\).

\begin{corollary}[Conditional Hilbert--P\'olya consequence]
\label{cor:hp}
In addition to Theorem~\ref{thm:deninger}, suppose \(H^1_{\mathbb C}\) has a
Hilbert completion \(\mathscr H^1\) on which \(\Theta_1\) has a densely
defined closed extension with compact resolvent.  Assume, with algebraic
multiplicities, the exact multiset equality
\[
 \operatorname{Spec}(\Theta_1)
 =\{\rho:\xi_{\mathbb R}(\rho)=0\},
 \qquad
 m_{\mathrm{alg}}(\rho,\Theta_1)
 =\operatorname{ord}_{s=\rho}\xi_{\mathbb R}(s),
\]
and assume that there is no additional point, continuous, or residual
spectrum.  Then
\[
 \operatorname{Spec}(\Theta_1)\subset\tfrac12+i\mathbb R
\]
is equivalent to the Riemann Hypothesis.

If, for a positive-definite Hilbert inner product,
\[
 \mathcal D(\Theta_1^*)=\mathcal D(\Theta_1),
 \qquad \Theta_1^*=1-\Theta_1,
\]
then \(H_{\mathrm{HP}}=-i(\Theta_1-\tfrac12)\) is self-adjoint.  Under the
same exact spectral comparison, this adjoint identity implies the Riemann
Hypothesis.
\end{corollary}

\begin{proof}
The first equivalence is the stated equality of multisets rewritten in
spectral language.  For the last assertion, put
\(A=\Theta_1-\tfrac12\).  The domain equality and adjoint identity say
\(A^*=-A\), so \(-iA\) is self-adjoint and
\(\operatorname{Spec}(A)\subset i\mathbb R\).
\end{proof}

For an unbounded operator, a formal adjoint formula checked only on a common
dense invariant domain proves only skew-symmetry unless essential
skew-adjointness is also shown. A compatible nondegenerate indefinite
pairing can yield reflected spectral symmetry without critical-line
confinement. In multiplicative language the eigenvalues
of \(e^{\Theta_1}\), not of \(\Theta_1\), would have modulus
\(e^{1/2}\) under the conclusion of Corollary~\ref{cor:hp}.
